\chapter[Aprendizados e Perspectivas]{Aprendizados e Perspectivas}
\label{resultados}

Este capítulo apresenta a análise crítica dos resultados obtidos a partir da implementação da solução. São discutidos os dados coletados nos testes, as evidências qualitativas, comparações com abordagens similares e as limitações técnicas observadas ao longo do desenvolvimento.

\section{Resultados Quantitativos e Qualitativos}

Os resultados quantitativos foram obtidos a partir dos testes de desempenho, confiabilidade e usabilidade, conforme descritos anteriormente. A Tabela~\ref{tab:resultados_quantitativos} resume os principais indicadores obtidos:

\begin{table}[H]
\centering
\caption{Indicadores de desempenho da solução}
\label{tab:resultados_quantitativos}
\begin{tabularx}{\textwidth}{|l|X|}
\hline
\textbf{Métrica} & \textbf{Valor Obtido} \\
\hline
Tempo médio de resposta & 180 ms \\
\hline
Taxa de sucesso das requisições & 99,3\% \\
\hline
Cobertura de testes unitários & 85\% \\
\hline
Disponibilidade durante testes & 100\% (ambiente de homologação) \\
\hline
\end{tabularx}
\end{table}

Além disso, foram obtidos feedbacks qualitativos sobre a usabilidade da aplicação por meio de simulações com usuários e análise heurística. Os pontos positivos mais citados foram a simplicidade da interface e a clareza dos relatórios gerados.

\section{Comparações com Soluções Existentes}

A solução proposta foi comparada com abordagens similares disponíveis no mercado ou em sistemas previamente utilizados pela organização. Em relação às soluções existentes, destacam-se os seguintes diferenciais:

\begin{itemize}
  \item Maior nível de automação e integração com sistemas legados;
  \item Interface responsiva e adaptada a múltiplos dispositivos;
  \item Maior cobertura de funcionalidades em um único sistema consolidado;
  \item Facilidade de manutenção e expansão futura.
\end{itemize}

\section{Interpretação dos Dados}

A análise dos dados obtidos indica que a solução atende de forma satisfatória os requisitos funcionais e não funcionais propostos. O tempo de resposta médio ficou abaixo do limite especificado (300 ms), e a taxa de sucesso de requisições demonstra estabilidade do sistema.

A cobertura dos testes também evidencia um bom nível de confiabilidade na lógica implementada. A aceitação positiva durante as simulações de uso reforça a aplicabilidade prática da ferramenta no ambiente-alvo.

\section{Limitações Observadas}

Apesar dos resultados positivos, algumas limitações técnicas e operacionais foram identificadas:

\begin{itemize}
  \item Dependência de conectividade estável com sistemas externos, o que pode impactar o desempenho em ambientes com infraestrutura limitada;
  \item Ausência de testes em ambiente de produção real com carga de usuários simultâneos elevada;
  \item Algumas funcionalidades foram simplificadas para garantir o escopo viável dentro do tempo do projeto;
  \item A interface administrativa ainda requer melhorias em termos de acessibilidade e personalização.
\end{itemize}

Essas limitações não comprometem o uso da solução, mas devem ser consideradas em futuras evoluções ou implantações em larga escala.

